\section{Symptomatic vs.~causal assessment}
\label{section:related-work:symptoms}


\begin{readerguide}
 It might be better to skip this section and to read it after the first
 high-level assessment is completed.
\end{readerguide}

\begin{wip}
 This part is yet to be written and will basically point out how the POP methodology differs from what we do here.
\end{wip}


Performance assessment can be guided by the observed behaviour or by monitoring the machine.
One might refer to this as a symptomatic vs.~causal assessment. 
Scientists will always favour a causal assessment where a hard metric can be used to explain certain observed behaviour.
However, this is not always possible and frequently not the most economic way to assess a code's performance.
Notably in early phases of a performance analysis, it makes sense to focus on what is observed without digging into the reasons for some behaviour.
